\section{引言：从 AlphaGo 到 AlphaProof}
Thomas Hubert 把这场讲座定位在 Berkeley LLM Agents SP25 的“推理-执行”单元，他用 AlphaZero/AlphaTensor 把强化学习的黄金时代串联到 AlphaProof。Hubert 明确指出，数学推理所需的\textbf{严密性}、\textbf{证明性 feedback} 与\textbf{无限创造力}，与 RL 边界探索所需的\textbf{搜索}、\textbf{泛化} 和\textbf{可靠回报}，是高度同构的。这就解释了为何 DeepMind 会把数学放在“RL 下一个战场”。

\begin{knowledgebox}{数学与 RL 的共鸣}
数学是一张高维网络，节点是定理，边是逻辑链；RL 则在动作空间做探索。Hubert 强调：“只要你能把定理抽象成 action policy，理论就能交给优化器去学习。”
\end{knowledgebox}

讲座分三段：第一段回溯数学形式化的演进并梳理 Lean/Mathlib；第二段总结 AlphaZero 的飞轮以及 AlphaProof 如何与证明系统对接；第三段聚焦工程实践、IMO 复现与治理闭环。整堂课按“理论→平台→应用”的教学逻辑展开，每段都留出时间展示经验教训与问答。

\begin{figure}[H]
\centering
\includegraphics[width=0.8\textwidth,height=0.35\textheight,keepaspectratio]{\videocoverpath}
\caption{Lecture 05 封面与 AlphaProof 流程示意（封面提供主视觉/架构意象，强调 RL + 形式化的融合）。}
\end{figure}

Hubert 特别提到，这堂课是 SP25 的第 5 讲，紧接在“生成式 agents 与 reasoning”之后。他把 AlphaProof 视作“数学版的 agent”，需要把 deterministic reward 和 interactive governance 结合，才能让高可信的 proof search 成为可能。

\subsection{本章小结}
介绍部分负责搭建语境：数学与 RL 之间的结构性相似性、AlphaProof 的使命、以及整个 SP25 讲座在“推理链 + 工程实践”之间的节奏定位，为后续内容奠定因果路径。
\newpage

